\chapter{Introdução}

\section{Visão Geral}

Nas esferas industrial e acadêmica de bancos de dados, a concepção de que um único modelo
de dados não atende a todas as necessidades tem conquistado aceitação, apontando para a
perspectiva de persistência poliglota no futuro. Embora os sistemas de banco de dados
relacionais devam manter sua predominância, espera-se uma convivência com diversas formas
de sistemas, como os NoSQL \cite{sadalage2013nosql}.

Algumas pesquisas têm sido feitas na área de persistência poliglota, como sugestões de
modelos de dados unificados \cite{chillon2023propagating}, comparação de desempenho de
bancos de dados multi-modelos contra abordagens de persistência poliglota
\cite{ye2023benchmark}, revisões sistemáticas sobre modelagem poliglota de dados
\cite{silva2024modelagem}, tutoriais sobre o estado da arte e desafios abertos
\cite{kiehn2022polyglot}, e a adoção da persistência poliglota em Big Data
\cite{khine2019review}.

Dada essa necessidade crescente da adoção de persistência poliglota pelas aplicações,
percebe-se a falta de uma ferramenta única de importação simultânea de dados para diversos
SGBDs que tratem modelos de dados diferentes. Até onde se investigou (Capítulo 3), não
foram identificadas ferramentas equivalentes: os importadores nativos atendem a um único
SGBD por vez, e as soluções de migração ou de modelagem poliglota não realizam a
materialização simultânea de um mesmo CSV em múltiplos paradigmas de dados --- o que
ressalta a originalidade desta proposta.

Assim sendo, a intenção deste trabalho é desenvolver a ferramenta \textit{Polyglot Import
CSV}, que importa dados armazenados em CSV para SGBDs relacionais e NoSQL. A configuração
da ferramenta será baseada em dois arquivos JSON: um de importação, mapeando entidades e
relacionamentos em cada SGBD, considerado seu modelo de dados; outro de configuração de
conexão dos SGBDs envolvidos. A escolha pelo formato JSON deve-se à sua simplicidade e
ampla adoção; sua finalidade é detalhar, de forma declarativa, o mapeamento das colunas do
CSV para os modelos de dados de cada SGBD, permitindo a validação prévia da importação.

CSV (\textit{Comma-Separated Values}) é um arquivo de valores separados por vírgulas,
frequentemente visualizado em Excel e outras planilhas eletrônicas. Pode haver outros tipos
de valores como delimitadores, mas o mais comum é a vírgula. Muitos sistemas e processos
hoje convertem seus dados para o formato CSV para saídas de arquivo para outros sistemas,
relatórios amigáveis para humanos, disponibilização de dados públicos, e outras
necessidades. É um formato de arquivo padrão com o qual humanos e sistemas já estão
familiarizados ao usar e manipular.

\section{Objetivos}

\subsection{Objetivo Geral}

Desenvolver a ferramenta \textit{Polyglot Import CSV} que possibilitará a importação de
dados em arquivos CSV para múltiplos SGBDs: PostgreSQL, Redis, MongoDB, Cassandra ou
Neo4j.

\subsection{Objetivos Específicos}

\begin{itemize}
  \item Projetar os arquivos de configuração JSON para que eles sigam sintaxes específicas,
        e flexíveis para lidar com vários modelos de dados, além de adicionar opções como
        filtros, escolha de chave primária e escolha do nome da coluna no banco de dados;
  \item Não permitir configurações incorretas, ou seja, se houver erro em algum lugar dos
        arquivos de configuração, a importação não deve começar até que sejam corrigidos
        (validação prévia da configuração);
  \item Avaliar o desempenho da ferramenta para diferentes volumes e modelos de dados
        destino;
  \item Disponibilizar a ferramenta de forma acessível ao usuário: inicialmente por meio de
        uma interface de linha de comando (CLI), com possibilidade de evolução para uma
        interface gráfica em trabalhos futuros.
\end{itemize}

\section{Metodologia}

O desenvolvimento deste trabalho segue uma metodologia de pesquisa aplicada, adotando as
seguintes etapas:

\begin{enumerate}
  \item Estudo sobre persistência poliglota e seu estado da arte;
  \item Estudo sobre os paradigmas não-relacionais que serão abrangidos pela solução:
        chave-valor, documento, grafo e colunar;
  \item Estudo de decisão de linguagem de implementação que demonstra melhor produtividade
        para codificar a solução;
  \item Projeto dos arquivos de configuração JSON;
  \item Projeto da instrução principal de importação que utiliza os arquivos de
        configuração;
  \item Implementação da solução ``Polyglot Import CSV'';
  \item Testes da solução ``Polyglot Import CSV'' implementada.
\end{enumerate}

Na primeira etapa, foi estudado o que é a persistência poliglota, e a sua relevância no
cenário atual.

Na segunda etapa, foram estudados os diversos modelos de dados não-relacio-
nais: como cada
modelo representa suas entidades; se há relacionamentos entre entidades no modelo; se é
possível a existência de entidades aninhadas e multivaloradas; as limitações de cada
modelo.

Na terceira etapa, foram estudadas as formas de implementar a solução usando as linguagens
Java e Python e, pela simplicidade e produtividade, a linguagem Python foi escolhida.

Na quarta e quinta etapas, foram projetados os dois arquivos de configuração JSON --- um
para o mapeamento da importação (entidades, relacionamentos e colunas) e outro para a
conexão dos SGBDs --- bem como a instrução principal de importação que os utiliza.

As últimas etapas, que estão em andamento, tratam da implementação da solução proposta,
``Polyglot Import CSV'', e dos testes da mesma com diferentes volumes de dados CSV e
importações simples e complexas para diferentes modelos de dados destino.

\section{Estrutura do Trabalho}

Este trabalho está estruturado em cinco capítulos principais. O primeiro aborda os
objetivos, a metodologia e o contexto que levaram à elaboração deste TCC. O capítulo 2
revisa fundamentos teóricos. O capítulo 3 apresenta trabalhos relacionados. O capítulo 4
descreve a proposta da ferramenta \textit{Polyglot Import CSV} e o capítulo 5 as
atividades futuras (TCC II).
